\sect{ऋषिपञ्चमी-व्रत-कथा}
\centerline{\small{(मूलम्—श्रीभविष्योत्तरपुराणम्, यथा श्रीव्रतरत्नाकरे संगृहीतम्)}}

\twolineshloka
{श्रुतानि देवदेवेश व्रतानि सुबहूनि च}
{साम्प्रतं मे समाचक्ष्व व्रतं पापप्रणाशनम्} %॥१॥

\uvacha{ब्रह्मोवाच}

\twolineshloka
{शृणु राजन् प्रवक्ष्यामि व्रतानामुत्तमं व्रतम्}
{ऋषिपञ्चमीति विख्यातं सर्वपापहरं परम्} %॥२॥

\twolineshloka
{येन चीर्णेन राजेन्द्र नरकं नैव पश्यति}
{अत्रैवोदाहरन्तीममितिहासं पुरातनम्} %॥३॥

\twolineshloka
{वैदर्भे च द्विजवर उदको नाम नामतः}
{तस्य भार्या सुशीलेति पतिव्रतपरायणा} %॥४॥

\twolineshloka
{तस्या अपत्ययुगलं पुत्रो हि सुविभीषणः}
{अधीतवान् सुतस्तस्य वेदान् साङ्गपदक्रमान्} %॥५॥

\twolineshloka
{समानेन कुलीनेन सुता चापि विवाहिता}
{नवोढात्वेऽपि सा दैवाद्वैधव्यं प्राप सत्तमा} %॥६॥

\twolineshloka
{सतीत्वं पालयन्ती सा आस्ते निजपितुर्गृहे}
{तस्या दुःखेन सन्तप्तस्सुतं संस्थाप्य वेश्मनि} %॥७॥

\twolineshloka
{गङ्गातीरवनं प्राप्तः कदाचित्तु तया सह}
{स तत्राध्यापयामास शिष्यान् वेदं द्विजोत्तमः} %॥८॥

\twolineshloka
{सुता च कुरुते तस्य पितुश्शुश्रूषणं परम्}
{पितुश्शुश्रूषणं कृत्वा परिश्रान्ता कदाचन} %॥९॥

\twolineshloka
{निशीथे किल सुप्तायां क्रिमिराशिरजायत}
{शिष्यास्तथाविधां दृष्ट्वा विवस्त्रां प्रस्तरे स्थिताम्} %॥१०॥

\twolineshloka
{गत्वा निवेदयामासुस्तन्मातुः करुणां गताः}
{न जानीमो वयं किञ्चिद्देवि साध्वी तथाविधा} %॥११॥

\twolineshloka
{क्रिमिराशिमयी जाता मातस्सम्प्रति दृश्यते}
{वज्रपातसमं श्रुत्वा तच्छिष्यैस्समुदाहृतम्} %॥१२॥

\twolineshloka
{सा श्रान्तमानसा शीघ्रं तत्समीपमुपागमत्}
{सुतां तथाविधां दृष्ट्वा विललाप सुदुःखिता} %॥१३॥

\twolineshloka
{उरश्च ताडयामास सुतरां मोहमाप च}
{क्षणेन प्राप्य चैतन्यमुत्थाप्य परिमृज्य च} %॥१४॥

\twolineshloka
{समालिङ्ग्य च बाहुभ्यां निन्ये तत्पितुरन्तिकम्}
{स्वामिन् कथय मे साध्वी केन दुष्कृतकर्मणा} %॥१५॥

\twolineshloka
{निशीथे च प्रसुप्तेयं जायते क्रिमिसङ्कुला}
{एवं श्रुत्वा ततो वाक्यम् ऋषिर्ध्यानपरायणः} %॥१६॥

\onelineshloka
{ज्ञात्वा निवेदयामास तस्याः प्राग्जन्मचेष्टितम्} %॥१७॥

\uvacha{ऋषिरुवाच}

\twolineshloka
{प्राणीयं सप्तमेऽतीते जन्मनि ब्राह्मणी ह्यभूत्}
{रजस्वला च सञ्जाता भाण्डादीनस्पृशत्तदा} %॥१८॥

\twolineshloka
{अस्यास्तु पाप्मना तेन जायते क्रिमिवद्वपुः}
{रजस्वलायाः पापेन युक्ता भवति साऽनघे} %॥१९॥

\twolineshloka
{प्रथमेऽहनि चण्डाली द्वितीये ब्रह्मघातिनी}
{तृतीये रजकी प्रोक्ता चतुर्थेऽहनि शुद्ध्यति} %॥२०॥

\twolineshloka
{तदा तया सखीसङ्घातं दृष्ट्वाऽवमानितम्}
{दृष्टव्रतप्रभावेन जाता द्विजकुलेऽमले} %॥२१॥

\twolineshloka
{अवमानाद्व्रतस्यास्य क्रिमिराशिरिवाधुना}
{एतत्ते कथितं सर्वं कारणं दुष्कृतस्य च} %॥२२॥

\uvacha{सुशीलोवाच}

\twolineshloka
{दर्शनादपि जन्तोरस्य द्विजानां निर्मले कुले}
{जन्म युष्मद्विधानां हि जायते ब्रह्मतेजसाम्} %॥२३॥

\twolineshloka
{अवज्ञया प्रजायन्ते निशीथे क्रिमिराशयः}
{महाश्चर्यं कथं नाथ तद्व्रतं कथयस्व मे} %॥२४॥

\uvacha{उदक उवाच}

\twolineshloka
{सुशीले शृणु तत्सम्यक् व्रतं नामोत्तमं व्रतम्}
{येन चीर्णेन सहसा पापादस्मात् प्रमुच्यते} %॥२५॥

\twolineshloka
{दुःखत्रयाच्च मुच्येत नारी सौभाग्यमाप्नुयात्}
{कल्याणानि विवर्धन्ते सम्पदश्च निरापदः} %॥२६॥

\uvacha{युधिष्ठिर उवाच}

\twolineshloka
{श्रुतानि देवदेवेश व्रतानि सुबहूनि च}
{साम्प्रतं त्वन्यदाचक्ष्व व्रतं पापप्रणाशनम्} %॥२७॥

\uvacha{श्रीकृष्ण उवाच}

\twolineshloka
{अथान्यदपि राजेन्द्र पञ्चम्यामृषिपूजनम्}
{कथयिष्यामि यत्कृत्वा नारी पापात् प्रमुच्यते} %॥२८॥

\uvacha{युधिष्ठिर उवाच}

\twolineshloka
{कीदृशी पञ्चमी कृष्ण कथं तदृषिपूजनम्}
{पातकान्मुच्यते कस्मान्नारी यदुकुलोद्भव} %॥२९॥

\twolineshloka
{पापानि च बहून्यत्र विद्यन्ते किल केशव}
{कथं वा ऋषिपञ्चम्यां नारी कस्मात् प्रमुच्यते} %॥३०॥

\uvacha{श्रीकृष्ण उवाच}

\twolineshloka
{अज्ञानाज्ज्ञानतो वापि या स्त्री जाता रजस्वला}
{दुष्टा स्पृशति भाण्डानि गृहकर्मणि संस्थिता} %॥३१॥

\twolineshloka
{प्राप्नोति सा महापापं सर्वा ना रजस्वलाः}
{शृणु तत्कारणं येन वर्जनीया रजस्वला} %॥३२॥

\twolineshloka
{प्रोत्सार्या गृहतो दूरं चातुर्वर्ण्येन भारत}
{ब्रह्महत्यां पुरा शक्रो वृत्रं हत्वा ह्यवाप च} %॥३३॥

\twolineshloka
{तया वै राजशार्दूल पीडितो वृत्रसूदनः}
{ब्रह्माणं समुपागच्छदात्मशुद्धिकारणात्} %॥३४॥

\twolineshloka
{शुद्ध्यै शक्रस्य राजेन्द्र प्रहृष्टेनान्तरात्मना}
{विभज्य ब्रह्महत्यां तु चतुर्धा स चतुर्मुखः} %॥३५॥

\twolineshloka
{प्राक्षिपद्राजशार्दूल चतुःस्थानेषु वै तदा}
{वह्नौ प्रथमजिह्वासु नदीनां प्रथमोदके} %॥३६॥

\twolineshloka
{गिरिवृक्षादिभूमौ च नारीरजसि पार्थिव}
{ततो रजस्वला नारी प्रोत्सार्या च प्रयत्नतः} %॥३७॥

\twolineshloka
{ब्राह्मणशासनात्पार्थ चातुर्वर्ण्येन सर्वदा}
{प्रथमेऽहनि चण्डाली द्वितीये ब्रह्मघातिनी} %॥३८॥

\twolineshloka
{तृतीये रजकी प्रोक्ता चतुर्थेऽहनि शुद्ध्यति}
{अज्ञानाज्ज्ञानतो वापि जातं सम्पर्कपातकम्} %॥३९॥

\twolineshloka
{तत्पापसंक्षयार्थं वै कार्येयमृषिपञ्चमी}
{सर्वपापप्रशमनी सर्वोपद्रवनाशिनी} %॥४०॥

\onelineshloka
{ब्रह्मक्षत्रियविट्शूद्रस्त्रीभिः कार्या विशेषतः} %॥४१॥

\uvacha{कृष्ण उवाच}

\onelineshloka
{अत्रार्थे यत्पुरा वृत्तं प्रवक्ष्यामि कथात्मकम्} %॥४२॥

\twolineshloka
{पुरा कृतयुगे राजा विदर्भाणां बभूव ह}
{श्येनजिन्नाम राजर्षिश्चातुर्वर्ण्यानुपालकः} %॥४३॥

\twolineshloka
{तस्य देशे वसन्विप्रो वेदवेदाङ्गपारगः}
{सुमित्रो नाम विप्रेन्द्रः सर्वभूतहिते रतः} %॥४४॥

\twolineshloka
{कृषिवृत्त्या सदा युक्तः कुटुम्बपरिपालकः}
{तस्य भार्या सुसाध्वी च पतिशुश्रूषणे रता} %॥४५॥

\twolineshloka
{जयश्रीर्नाम विख्याता बहुभृत्यसुहृज्जना}
{अतिचिन्तान्विता सा च प्रावृट्काले सुमध्यमा} %॥४६॥

\twolineshloka
{क्षेत्रादिषु रता साध्वी चाकुलीकृतमानसा}
{एकदा साऽऽत्मनः प्राप्तमृतुकालं व्यलोकयत्} %॥४७॥

\twolineshloka
{रजस्वलाऽपि सा राजन् गृहकर्म चकार ह}
{भाण्डादीन्यस्पृशद्राजन् प्राप्तर्तुरपि भामिनी} %॥४८॥

\twolineshloka
{कालेन बहुना साध्वी पञ्चत्वमगमत्तदा}
{तस्या भर्ताऽपि विप्रोऽसौ कालधर्ममुपेयिवान्} %॥४९॥

\twolineshloka
{एवं तौ दम्पती राजन् स्वकर्मवशगौ तदा}
{भार्या तस्य जयश्रीस्सा ऋतुसम्पर्कदोषतः} %॥५०॥

\twolineshloka
{शुनीयोनिसमुत्पन्ना शुनी चाभून्नरेश्वर}
{तस्यास्सम्पर्कदोषेण बलीवर्दो बभूव सः} %॥५१॥

\twolineshloka
{ऋतुसम्पर्कदोषेण तिर्यग्योनिमुपागतौ}
{स्वधर्माचरणाज्जातावुभौ जातिस्मरौ तदा} %॥५२॥

\twolineshloka
{सुतस्यैव गृहे राजन् स्मरन्तौ पूर्वपातकम्}
{सुमित्रस्य च पुत्रोऽभूद्गुरुशुश्रूषणे रतः} %॥५३॥

\twolineshloka
{सुमतिर्नाम धर्मज्ञो देवतातिथिपूजकः}
{कालक्षयाहे सम्प्राप्ते पितुस्तु सुमतिस्तदा} %॥५४॥

\twolineshloka
{भार्यां चन्द्रवतीं प्राप्य पैतृकश्रद्धयाऽन्वितः}
{अद्य साम्वत्सरदिनं पितुर्मे चारुहासिनि} %॥५५॥

\twolineshloka
{भोजनीया द्विजा भीरु पाकशुद्धिर्विधीयताम्}
{तया कृता पाकशुद्धिस्सुमते शासनाद्गृहे} %॥५६॥

\twolineshloka
{मुक्तं पायसभाण्डे वै सर्पेण गरलं ततः}
{दृष्ट्वा ब्रह्मवधाद्भीता शुनी भाण्डानि साऽस्पृशत्} %॥५७॥

\twolineshloka
{द्विजभार्या च तां दृष्ट्वा उल्मुकेन जघान ह}
{भाण्डादीनि च प्रक्षाल्य त्यक्त्वा पाकं सुमध्यमा} %॥५८॥

\twolineshloka
{पुनः पाकं च कृत्वा तु श्राद्धं कृत्वा विधानतः}
{ततो भुक्तेषु विप्रेषु नोच्छिष्टं च ददौ बहिः} %॥५९॥

\twolineshloka
{भूमौ क्षिप्तं तया शुन्या उपवासस्ततोऽभवत्}
{ततो रात्र्यां प्रवृत्तायां सा शुनी क्षुधिता भृशम्} %॥६०॥

\onelineshloka
{बलीवर्दमुपागत्य भर्तारमिदमब्रवीत्} %॥६१॥

\uvacha{शुन्युवाच}

\twolineshloka
{बुभुक्षिता गृहे भर्तर्न दत्तं भोजनादिकम्}
{ग्रासादिकं च न प्राप्तं क्षुधा मां बाधते भृशम्} %॥६२॥

\twolineshloka
{अन्यस्मिन् दिवसे पुत्रो मम लेह्यं ददात्यसौ}
{अद्य मह्यं किमप्येष उच्छिष्टमपि नो ददौ} %॥६३॥

\twolineshloka
{पायसान्ने पपाताद्य गरलं सर्पसम्भवम्}
{मया विचिन्त्य मनसा मरिष्यन्ति द्विजोत्तमाः} %॥६४॥

\twolineshloka
{संस्पृष्टं पायसं गत्वा तत्राहं ताडिता भृशम्}
{दुःखितं तेन मे गात्रं कटिर्भग्ना करोमि किम्} %॥६५॥

\uvacha{बलीवर्द उवाच}

\onelineshloka
{भद्रे ते पापजं फलम्} %॥६६॥

\twolineshloka
{किं करोमि शक्तोऽहं भारवाहोऽद्य सुस्थितः}
{अद्याहमात्मजक्षेत्रे वाहितस्सकलं दिनम्} %॥६७॥

\twolineshloka
{मारितश्चात्मजेनाहं मुखं बद्ध्वा बुभुक्षितः}
{वृथा श्राद्धं कृतं तेन जाताऽद्य मम कष्टता} %॥६८॥

\uvacha{श्रीकृष्ण उवाच}

\twolineshloka
{तयोस्संवदतोरेव मातापित्रोश्च भारत}
{श्रुत्वा पुत्रस्तदा वाक्यं यदुक्तं च तदोभयोः} %॥६९॥

\twolineshloka
{ततो रजन्यां तत्कालं ददौ तस्यैव भोजनम्}
{पितरौ तौ विदित्वा तु दत्तवान् सुमतिस्तदा} %॥७०॥

\twolineshloka
{ततोऽसौ दुःखितः पुत्रो ज्ञात्वाऽवस्थां तदा तयोः}
{मातापित्रोस्तु राजेन्द्र दुःखात्मा प्रस्थितो वनम्} %॥७१॥

\twolineshloka
{ज्ञातुमिच्छामि कष्टं वै मातापित्रोश्च भारत}
{तत्र गत्वा ज्ञानवृद्धानृषीन् सप्त वनाश्रितान्} %॥७२॥

\onelineshloka
{प्रणिपत्याब्रवीद्वाक्यं हितं चैव तदा तयोः} %॥७३॥

\uvacha{सुमतिरुवाच}

\twolineshloka
{कथयध्वं विप्रवर्याः प्रश्नमेकं समाहिताः}
{केन कर्मविपाकेन पितरौ मे तपोधनाः} %॥७४॥

\onelineshloka
{इमामवस्थां सम्प्राप्तौ मुच्येते पातकात्कथम्} %॥७५॥

\uvacha{श्रीकृष्ण उवाच}

\twolineshloka
{तदाकर्ण्य वचस्तस्य सुमतेर्दुःखितस्य च}
{ऋषिस्सर्वतपा नाम सर्वज्ञः करुणान्वितः} %॥७६॥

\onelineshloka
{सुमतिं प्रत्युवाचेदं तत्पित्रोर्मुक्तये तदा} %॥७७॥

\uvacha{ऋषिरुवाच}

\onelineshloka
{तव माता पुरा विप्र स्वगृहे बालभावतः} %॥७८॥

\twolineshloka
{ऋतुं प्राप्तं विदित्वा तु सम्पर्कमकरोद्द्विज}
{तेन कर्मविपाकेन शुनीयोनिमुपागता} %॥७९॥

\twolineshloka
{पिताऽपि स्पर्शदोषेण बलीवर्दो बभूव ह}
{एतयोर्मुक्तिकामार्थं कुरु त्वमृषिपञ्चमीम्} %॥८०॥

\twolineshloka
{भार्यया सह विप्रेन्द्र ऋषीन् सम्पूज्य यत्नतः}
{आचरस्व व्रतं तत्र सप्तवर्षं द्विजोत्तम} %॥८१॥

\twolineshloka
{अन्ते चोद्यापनं कुर्याद्वित्तशाठ्यविवर्जितः}
{शाकाहारस्तु कर्तव्यो नीवारैश्श्यामकैस्तदा} %॥८२॥

\twolineshloka
{कन्दं वाऽथ फलं मूलं हलकृष्टं न भक्षयेत्}
{प्राप्य भाद्रपदे मासि शुक्लपक्षस्य पञ्चमीम्} %॥८३॥

\twolineshloka
{तस्यां मध्याह्नसमये नद्यादौ विमले जले}
{कृत्वाऽपामार्गसमिधा दन्तधावनमादितः} %॥८४॥

\twolineshloka
{आयुर्बलं यशो वर्चः प्रजाः पशुवसूनि च}
{ब्रह्म प्रज्ञां च मेधां च त्वन्नो देहि वनस्पते} %॥८५॥

\twolineshloka
{संप्रार्थ्यानेन मन्त्रेण कुर्याद्वै दन्तधावनम्}
{मुखदुर्गन्धनाशाय दन्तानां च विशुद्धये} %॥८६॥

\twolineshloka
{ष्ठीवनाय च गात्राणां कुर्वेऽहं दन्तधावनम्}
{अनेन दन्तान् संशोध्य स्नायान्मृत्स्नानपूर्वकम्} %॥८७॥

\twolineshloka
{तिलामलककल्केन केशान् संशोध्य यत्नतः}
{परिधाय नवं शुद्धं वासोयुग्मं समाहितः} %॥८८॥

\twolineshloka
{विधाय नित्यकर्माणि गृहं गत्वा तान् ऋषीन्}
{स्थापयेद्विधिवद्भक्त्या पञ्चामृतरसैश्शुभैः} %॥८९॥

\twolineshloka
{चन्दनागरुकर्पूरैर्विलिप्य च सुगन्धिभिः}
{पूजयेद्विविधैः पुष्पैर्गन्धधूपादिदीपकैः} %॥९०॥

\twolineshloka
{समाच्छाद्य शुभैर्वस्त्रैस्सोपवीतैर्यथाविधि}
{ततो नैवेद्यसम्पन्नैर्दद्याच्छुभैः फलैः} %॥९१॥

\onelineshloka
{पूजयस्व ऋषीन् दिव्यान् अरुन्धत्या समन्वितान्} %॥९२॥

\twolineshloka
{कश्यपोऽत्रिर्भरद्वाजो विश्वामित्रोऽथ गौतमः}
{जमदग्निर्वसिष्ठश्च साध्वी चैवाप्यरुन्धती} %॥९३॥

\twolineshloka
{मन्त्रेणानेन सप्तर्षीन् पूजयेत्सुसमाहितः}
{व्रतेन ऋषिपञ्चम्याः कृतेनैव द्विजोत्तम} %॥९४॥

\onelineshloka
{ऋतुसम्पर्कजो दोषः क्षयं याति न संशयः} %॥९५॥

\uvacha{श्रीकृष्ण उवाच}

\onelineshloka
{तच्छ्रुत्वा सुमतिर्वाक्यं परमं मुनिभाषितम्} %॥९६॥

\twolineshloka
{गृहमेत्य व्रतं चक्रे सभार्यः श्रद्धयाऽन्वितः}
{व्रतं तु ऋषिपञ्चम्यास्सर्वपापप्रणाशनम्} %॥९७॥

\twolineshloka
{कृत्वा सर्वं यथोक्तं च मातापित्रोः फलं ददौ}
{व्रतपुण्यप्रभावेन माता तस्य श्वयोनितः} %॥९८॥

\twolineshloka
{मुक्ता नृपतिशार्दूल विमानवरमास्थिता}
{दिव्याम्बरधरा भूत्वा गता स्वर्गं च भारत} %॥९९॥

\twolineshloka
{पिताऽपि स मृतोऽगच्छत् सुमतेः पशुयोनितः}
{स्वर्गं प्राप्तो महाराज व्रतस्यास्य प्रभावतः} %॥१००॥

\twolineshloka
{दिव्याम्बरधरो भूत्वा गतस्स्वर्गं स भारत}
{कायिकं वाचिकं वापि मानसं यच्च दुष्कृतम्} %॥१०१॥

\twolineshloka
{तत्सर्वं विलयं याति व्रतस्यास्य प्रभावतः}
{तस्य यज्जायते पुण्यं तच्छृणुष्व नृपोत्तम} %॥१०२॥

\twolineshloka
{सर्वव्रतेषु यत्पुण्यं सर्वतीर्थेषु यत्फलम्}
{सर्वदानेषु दत्तेषु तदेतद्व्रतचारणात्} %॥१०३॥

\twolineshloka
{कुरुते या व्रतं नारी सा भवेत्सुखभागिनी}
{रूपलावण्ययुक्ता च पुत्रपौत्राभिसंयुता} %॥१०४॥

\twolineshloka
{इह लोके सुखं यायात् परत्र च परां गतिम्}
{एतत्ते कथितं राजन् व्रतानामुत्तमं व्रतम्} %॥१०५॥

\twolineshloka
{सर्वसम्पत्प्रदं चैव नारीणां पापनाशनम्}
{धनं यशश्च स्वर्गं च पुत्रानपि युधिष्ठिर} %॥१०६॥

\onelineshloka
{पठतां शृण्वतां चापि दद्यात्पापप्रणाशनम्} %॥१०७॥

॥ इति श्रीभविष्योत्तरपुराणे ऋषिपञ्चमीव्रतकथा सम्पूर्णा ॥
